\subsection{Cross-Domain Span Completion}
\label{sec:cross_domain_completion}

\paragraph{Task Description.}
The cross-domain completion experiment evaluates the substrate's
domain-agnosticism: without any domain-specific ingestion, indexing,
or parameter adjustment, the same value manifold is asked to complete
the final 50 bytes of samples drawn from three structurally distinct
domains:

\begin{itemize}[nosep]
  \item \textbf{Natural Language} --- English Wikipedia
        (\texttt{wikimedia/wikipedia}, subset \texttt{20231101.en})
  \item \textbf{Source Code} --- Python source files
        (\texttt{bigcode/the-stack-smol})
  \item \textbf{Biology} --- Amino acid + DSSP3 structure labels
        (\texttt{proteinea/secondary\_structure\_prediction})
\end{itemize}

Each domain contributes $100$ training samples,
ingested sequentially into a single unified substrate.
Test prompts hold out the last 50 bytes; the system reconstructs
them from the value resonance field without any domain indicator.

\paragraph{Results.}
Figure~\ref{fig:cross_domain_map} shows the composite.
Panel~A is the per-sample score fingerprint heatmap; sample labels
encode domain (first four characters) and local index.
Panel~B shows mean Exact / Partial / Fuzzy / Weighted scores grouped
by domain, directly comparing substrate performance across the
three data modalities.

Across all $N = 1$ samples the overall weighted score was
0.000. Per-domain weighted means:
\textbf{Natural Language}: 0.000 (exact: 0.0\%).

\paragraph{Assessment.}
Completion accuracy was low across all domains at this ingestion
scale.  The result is consistent with the theoretical expectation:
the attractor field requires a minimum density of related samples
before value resonance can reliably recover novel suffixes.

